\chapter{What Is CNA?}
\label{ch:what-is-cna}

\section{The one-sentence version}

CNA is a C\texttt{++}23 reimplementation of the \emph{XNA 4.0} programming model --- the
game framework Microsoft shipped for the Xbox 360 and Windows from roughly 2006 to 2013 ---
built on top of SDL3, with the actual rendering delegated to one of a dozen interchangeable
graphics backends chosen at build time.

That sentence hides three separate engineering problems, and it is worth pulling them
apart before going any further, because the rest of this book is organized around exactly
that separation:

\begin{enumerate}
  \item \textbf{Reproduce an API.} XNA's public surface --- \cnans{Microsoft.Xna.Framework}
        and its sub-namespaces --- has to exist in C\texttt{++}, with the same class names,
        method signatures, namespaces, and observable behavior, so that code written against
        real XNA (or against FNA, the open-source XNA reimplementation in C\# that this
        project treats as its primary behavioral reference) has somewhere to go.
  \item \textbf{Reproduce a runtime.} XNA is a managed-code API. It assumes things a bare
        C\texttt{++} program does not have for free: garbage-collected \texttt{System.Object},
        events and delegates, exceptions with a .NET-shaped hierarchy, \texttt{TimeSpan},
        collections. CNA does not reimplement the CLR --- instead it depends on a sibling
        project, \texttt{sharp-runtime}, which provides a deliberately partial C\texttt{++}
        subset of the .NET base class library, just deep enough to make the XNA API's
        vocabulary make sense in C\texttt{++}.
  \item \textbf{Reproduce a GPU.} Real XNA only ever had to run against Direct3D 9 (Xbox 360's
        own GPU API is Direct3D-shaped too). CNA has to make the same \cnaclass{GraphicsDevice}
        /\cnaclass{SpriteBatch}/\cnaclass{Effect} surface work identically whether the pixels
        are actually produced by SDL's own 2D renderer, raw OpenGL, Vulkan, BGFX, WebGPU, real
        Direct3D 9/11/12 under Wine, an HTML \texttt{<canvas>} under Emscripten, or a
        from-scratch DirectDraw reimplementation. That is the job of the backend layer covered
        in Part~IV of this volume.
\end{enumerate}

\begin{sourcenote}
\textbf{From the project's own README (2026-07-11 snapshot):} ``CNA is a C\texttt{++}
reimplementation of the XNA 4.0 programming model, built on SDL3 and a pluggable graphics
backend layer. It is a framework/runtime and abstraction layer---not a game---designed to
preserve XNA-style APIs (\texttt{Microsoft::Xna::Framework}) while using modern C\texttt{++}
internals.''
\end{sourcenote}

\section{Why XNA, and why now}

XNA Game Studio was Microsoft's managed-code (C\#) game development framework, built around
a small, coherent \cnaclass{Game}/\cnaclass{GraphicsDevice}/\cnaclass{SpriteBatch}/\cnaclass{Effect}
vocabulary that a whole generation of independent and student developers learned game
programming on. Microsoft discontinued active XNA development years ago; the community
response was \textbf{FNA}, a faithful open-source reimplementation of the XNA 4.0 API in
C\# on top of SDL2, which today is what a great many surviving XNA-era games (including
several Xbox Live Arcade and Steam titles) actually ship on.

CNA asks a different question than FNA does. FNA asks: \emph{how do we keep running C\#
XNA code forever, on modern platforms, without Microsoft?} CNA asks: \emph{what if the code
itself did not have to be C\#?} A native C\texttt{++} implementation of the same API gives a
team that likes the XNA/MonoGame mental model --- a thin, explicit, immediate-mode-ish 2D/3D
API without an editor or a scene graph imposed on you --- a path that does not require a
managed runtime or GC pauses, and that a build system and toolchain of your choosing can
own end to end. CNA's own README states this directly: it is meant to ``provide a native
C\texttt{++} path for teams that like the XNA/MonoGame model but need non-managed
runtime/toolchain control.''

This is also, not incidentally, why CNA is useful as a target for \emph{porting} existing
XNA/FNA games rather than only for writing new ones from scratch --- a use case this book
returns to in Volume II's migration and case-study chapters, using the real ports of
\emph{Speedy Blupi} and \emph{Planet Blupi} (via the \texttt{free-direct} sibling project)
as worked examples.

\section{What CNA is not}

It is worth being precise about the negative space, because CNA's own documentation is:

\begin{itemize}
  \item \textbf{Not a game.} There is no bundled playable demo executable shipped as the
        project's flagship artifact --- the README states this is a deliberate priority of
        framework/runtime development over a shipping game demo. What ships instead is a
        test suite (thousands of GoogleTest cases across the unit and GPU pixel-test suites)
        and a set of minimal verification executables like \texttt{hello-triangle-sdl}.
  \item \textbf{Not a full XNA content pipeline.} CNA deliberately does not implement a real
        \texttt{.xnb} binary content reader. Assets are loaded as raw files plus JSON
        descriptors instead. This is a named design choice, not an oversight --- see
        Chapter~\ref{ch:content-and-assets} and the migration guide discussion in Volume II.
  \item \textbf{Not compiled-shader compatible (yet).} The \cnaclass{Effect(GraphicsDevice\&,
        byte[])} constructor --- loading a precompiled \texttt{.fx} shader bytecode blob, the
        way a real XNA game's content pipeline output would --- is, in the project's own
        words, ``the single biggest real gap,'' blocking 23 of the 86 official XNA samples in
        the sibling \texttt{cna-samples} repository. Chapter~\ref{ch:shader-fx-gap} covers
        this in depth, including why it is architecturally hard rather than merely unwritten.
  \item \textbf{Not binary-compatible with real Xbox Live.} The \cnaclass{GamerServices}
        namespace (Chapter references in Volume II) reimplements the public API shape with
        local/synthetic semantics, matching how FNA itself already handles this namespace ---
        there is no real Xbox Live behind it, by design.
  \item \textbf{Not wire-compatible with Microsoft DirectPlay/DirectDraw/DirectSound.} The
        \texttt{free-direct} sibling project (Chapter~\ref{ch:dx3-freedirect-backend} and
        Volume II Part~III) that powers the \texttt{DX3} backend is explicit that it targets
        FreeDirect-to-FreeDirect compatibility for specific legacy games, never the real
        Microsoft wire protocol.
\end{itemize}

\section{A framework built to be audited, not just believed}

The single most distinctive thing about how this project is run --- and the reason this book
can make specific, checkable claims instead of vague ones --- is its verification
methodology. Three mechanisms recur throughout the codebase and throughout this book:

\begin{description}
  \item[Differential testing against a real reference implementation.] \texttt{tools/fna-reference/}
        runs a genuine, live \texttt{FNA.dll} and compares CNA's behavior against it
        call-for-call, rather than against a written specification of what FNA is supposed to
        do. Disputed behavior that differential testing cannot settle is checked against real
        XNA 4.0 running in a Windows~7 virtual machine.
  \item[A compile-time purity check.] The \texttt{NOXNA} build option turns every
        non-XNA-tagged declaration into a \texttt{[[deprecated]]} warning under
        \texttt{-Werror}. Because every CNA extension beyond the real XNA 4.0 API must be
        wrapped in the \texttt{NOXNA} marker macro (Chapter~\ref{ch:design-philosophy}
        explains the convention in full), this check gives the project a mechanical --- not
        just documented --- guarantee that the public \cnaclass{Microsoft::Xna::Framework}
        surface has not silently grown project-specific extensions.
  \item[Pixel-level oracle corpora, not just ``it compiles.''] The \texttt{D3D9} backend in
        particular ships a 31-scene oracle corpus that diffs CNA's own render output against
        \emph{the real XNA 4.0 runtime's own render} of the same scene, at zero tolerance. As
        of the writing that this book draws from, 0 of 31 scenes diverge. Chapter~\ref{ch:direct3d-backends}
        covers what that verification pipeline (MinGW-w64 cross-compilation, Wine, DXVK, a
        real GPU) actually looks like.
\end{description}

This matters for how you should read the rest of this book. Where a claim in the chapters
ahead is stated as a fact about CNA's behavior, it is because the project's own test suite,
documentation, or source comments state it as a verified fact, generally with a specific
mechanism (a CTest binary, a doc file, a task number) behind it --- and where the project's
own documentation flags something as a dated snapshot, an estimate, or a known open gap, this
book preserves that qualification rather than erasing it for a cleaner sentence.

\section{Reading map for the rest of this volume}

Chapter~\ref{ch:ecosystem-map} draws the full dependency graph between CNA and its sibling
repositories. Chapter~\ref{ch:design-philosophy} is the ``house style'' chapter --- the
concrete rules (API mirroring, namespace discipline, the \texttt{NOXNA} tag, the
\texttt{getXProperty}/\texttt{setXProperty} convention, the porting checklist) that every
other chapter's code examples silently assume. Part~II gets you from a clean checkout to a
running triangle and then through the core framework types. Part~III and Part~IV are the
graphics machine, first as an API and then as eleven different implementations of that same
API.
